\problemname{The Sawmill}

Alma has a lot of planks that she wants to saw into pieces. Therefore, she has gone to a state-of-the-art
sawmill that automatically cuts and sorts the planks. Each plank can be represented by an infinitely long
$x$-axis, and the sawmill will cut the plank at the $N+1$ positions $$x_1, x_2, \dots, x_N, v$$
where $v$ is a number chosen by the user. Then, all finitely long pieces are sorted and output by the sawmill.
Alma wants to know the numbers $x_1, x_2, \dots, x_N$, but it seems that no one knows what the inside of the
sawmill looks like. So instead, she plans to figure out the numbers by feeding in some planks with well-chosen values of $v$.

There are $N$ secret integers $x_1 < x_2 < \dots < x_N$ between $1$ and $10^9$. Note that all these numbers
are different. Your goal is to find these numbers. You can send planks to the sawmill. The sawmill takes
an integer $v$ as input ($1 \leq v \leq 10^9$), and does the following:

\begin{enumerate}
  \item Create the list $L = [x_1, x_2, \dots, x_N, v]$.
  \item Sort $L$.
  \item Create the list $D$, where $D_i = L_{i+1}-L_i$ for all $i = 1, 2, \dots, N$.
  \item Sort $D$, and return the $N$ integers in $D$.
\end{enumerate}

You may send at most $N$ planks, except for subtask 4, where you can send $N+1$ planks.

\section*{Interaction}
Your program should first read two integers $N$ and $T$ ($1 \leq N \leq 1000$, $1 \leq T \leq 5$).
$N$ is the number of secret numbers you need to find, and $T$ is the test group number.
The reason $T$ is given as input is to make it easier to earn partial points.

Then, you can start sending planks. Print a line with ``\verb|? v|'' to send a plank with the
number $v$ to the sawmill. The number $v$ must satisfy $1 \leq v \leq 10^9$. Then, your program
should read $N$ integers on one line, the numbers $D_1, D_2, \dots, D_N$. Note that $D_1$ can
be zero if $v$ was equal to one of the numbers $x_i$.

When you have found the numbers $x_1, x_2, \dots, x_N$, you should print a line in the form 
$$! \; x_1 \; x_2 \; x_3 \; \dots \; x_N$$
Then, your program should terminate and not print anything more.

\textbf{Make sure to flush the output after each query}, otherwise, you may get \textit{Time Limit Exceeded}.
In C++, this can be done with, for example, \texttt{cout << flush} or \texttt{fflush(stdout)};
in Python with \texttt{stdout.flush()};
and in Java with \texttt{System.out.flush()}.


\section*{Scoring}
Your solution will be tested on a set of test groups, each worth a number of points. Each test group contains
a set of test cases. To get the points for a test group you need to solve all test cases in the test group.

\noindent
\begin{tabular}{| l | l | p{12cm} |}
  \hline
  \textbf{Group} & \textbf{Points} & \textbf{Constraints} \\ \hline
  $1$   & $15$       & $N = 1$\\ \hline
  $2$   & $15$       & $N = 2$  \\ \hline
  $3$   & $11$       & $x_i \leq N+1$ \\ \hline
  $4$   & $37$       & $N \leq 100$, $x_i \leq 10^4$, you may send at most $N+1$ planks. \\ \hline
  $5$   & $22$       & No additional constraints. \\ \hline
\end{tabular}

